\chapter{Order Types and Their Uses}\label{ch:mx:order-types-and-their-uses}

In January 2015 two US stock exchanges agreed to pay a USD~14 million penalty, without admitting or denying the findings, to settle SEC charges
about their order types. Their rules described one way of re-pricing an order that would lock or cross another venue's quote; the exchanges
accepted three, with different priority against each other, and had described one of them completely to some members but not to all. Nothing was
hidden in the \omterm{def:mx:the-limit-order-book:engine}{matching engine}'s code. It was hidden in the difference between what the engine did with an order and what the rulebook said it
would do. An order type is exactly that: a contract about what the engine will do with an order, and its fine print decides who trades first.

\section{Price conditions: limit, market, marketable limit, stop}

A \omterm{def:mx:the-limit-order-book:orders}{limit order} and a \omterm{def:mx:the-limit-order-book:orders}{market order} (\cref{def:mx:the-limit-order-book:orders}) differ in what they guarantee: the \omterm{def:mx:the-limit-order-book:orders}{limit order} its price, the \omterm{def:mx:the-limit-order-book:orders}{market
order} its execution, as long as there is anything to execute against. Between them sits the order most traders actually send.

\begin{definition}[Marketable limit order]\label{def:mx:order-types-and-their-uses:marketable}
A \emph{marketable limit order}\index{marketable limit order} is a \omterm{def:mx:the-limit-order-book:orders}{limit order} whose price reaches the best opposite quote when it arrives: a buy at or
above the best ask, a sell at or below the best bid. It executes at once against the book, up to its limit price, and any remainder rests at that price.
\end{definition}

A \omterm{def:mx:order-types-and-their-uses:marketable}{marketable limit order} is a \omterm{def:mx:the-limit-order-book:orders}{market order} with a guard rail: a buyer who sends 200 shares at the ask plus five cents trades everything available up
to that price and no further. The guard matters when the book is thin, and it matters most when orders arrive that were not sent by a person
reacting to the market but by a price the market reached.

\begin{definition}[Stop order, stop-limit order]\label{def:mx:order-types-and-their-uses:stop}
A \emph{stop order}\index{stop order} rests unseen until a trade occurs at or through its stop price (at or above it for a buy stop, at or below it for
a sell stop); it then becomes a \omterm{def:mx:the-limit-order-book:orders}{market order}. A \emph{stop-limit order}\index{stop-limit order} becomes a \omterm{def:mx:the-limit-order-book:orders}{limit order} at its stated limit price instead.
\end{definition}

Stops are the order of traders who want out when a price is reached: the stop-loss. Their triggers are prices, their execution is immediate, and a
book full of them is a book that sells harder the more it falls. \Cref{fig:mx:order-types-and-their-uses:cascade} measures this on
\texttt{firm.exchsim}, the book's exchange simulator (built in chapter~26): a bid ladder of 500 shares at every cent below 100.00, sell stops of 800 shares
at every cent from 99.95 down, and one market sell of 3\,000 shares. Without stops the sale walks the ladder down five cents. With twenty stops beneath it,
each triggered sale takes out more bids, reaches the next trigger, and the price falls 37 cents; 19\,000 shares trade for an initial sale of 3\,000.
\omterm{def:mx:order-types-and-their-uses:stop}{Stop-limit orders} priced one cent below their triggers stop the cascade sooner (25 cents), because an order that cannot sell at its limit rests
instead of walking down: the guard rail again, bought at the price of not getting out.

\begin{omfigure}
\begin{tikzpicture}
\begin{axis}[width=8.6cm,height=5.2cm,xlabel={\small sell stops beneath the market},ylabel={\small price fall (cents)},
  tick label style={font=\footnotesize},xtick={0,2,4,6,8,10,12,14,16,18,20},ymin=0,ymax=40,grid=major,grid style={gray!25},
  legend columns=2,legend style={font=\footnotesize,at={(0.5,-0.3)},anchor=north,draw=none,fill=none,
  /tikz/every even column/.append style={column sep=8pt}},table/col sep=comma]
\addplot[omThm,very thick,mark=*] table[x=n,y=stop_market]{figdata/microstructure/02-order-types-and-their-uses/cascade.csv};
\addplot[omDef,very thick,mark=square*] table[x=n,y=stop_limit]{figdata/microstructure/02-order-types-and-their-uses/cascade.csv};
\legend{stop orders,stop-limit orders (one cent below the trigger)}
\end{axis}
\end{tikzpicture}

\omcaption{A stop cascade on \texttt{firm.exchsim}: the fall caused by one market sell of 3\,000 shares into a ladder of 500 shares a cent, against the
number of 800-share sell stops placed a cent apart from 99.95 down. Data: \texttt{mx\_ordertypes.stop\_cascade}.}\label{fig:mx:order-types-and-their-uses:cascade}
\end{omfigure}

\section{Time conditions}

\begin{definition}[Time in force; immediate-or-cancel, fill-or-kill and good-till-cancelled orders]\label{def:mx:order-types-and-their-uses:tif}
An order's \emph{time in force}\index{time in force} says how long its unexecuted part may rest. A day order rests until the end of the session. An
\emph{immediate-or-cancel order}\index{immediate-or-cancel order} (IOC) executes what it can on arrival and cancels the rest. A \emph{fill-or-kill
order}\index{fill-or-kill order} (FOK) executes in full on arrival or not at all. A \emph{good-till-cancelled order}\index{good-till-cancelled order}
(GTC) survives the session's end and rests until it is executed or cancelled.
\end{definition}

IOC and FOK orders never rest, so they never show anything to the market before they trade: they are how algorithms take liquidity, sweep several
venues at once (the intermarket sweep order of One Quant Book 1, chapter~9, is an IOC with a legal flag), and ping dark pools (chapter~9). The at-the-open
and at-the-close times in force, which rest only for an auction, belong to chapter~10. The simulator's order-entry protocol carries all six
(\texttt{D}, \texttt{G}, \texttt{I}, \texttt{F}, \texttt{O}, \texttt{C}).

The table shows what one book does with eight orders of different kinds. Firm~1 rests 300 shares bid at 99.99 and 500 offered at 100.01; firm~2 then
sends, in order:

\begin{center}
\small
\begin{tabular}{@{}p{4.4cm}p{3.2cm}rp{4.6cm}@{}}
\toprule
order & reports & filled & what happened \\
\midrule
limit buy 100 at 100.00 & accepted & 0 & rests; becomes the best bid, 100 displayed \\
marketable limit buy 200 at 100.01 & accepted, executed & 200 & takes 200 of the 500 offered \\
post-only buy 100 at 100.01 & rejected (would cross) & 0 & never enters the book \\
fill-or-kill buy 400 at 100.01 & accepted, cancelled & 0 & only 300 remain at 100.01: nothing trades \\
immediate-or-cancel buy 400 at 100.01 & accepted, executed, cancelled & 300 & takes the 300, cancels 100 \\
iceberg sell 1\,000 at 100.02, 100 shown & accepted & 0 & rests; the feed shows 100 \\
\omterm{def:mx:order-types-and-their-uses:peg}{midpoint peg} sell 200 & accepted & 0 & rests hidden at the mid, 100.01 \\
sell stop 300 at 99.99 & accepted (waiting) & 0 & invisible until a trade at 99.99 or lower \\
\bottomrule
\end{tabular}
\end{center}

\section{Display conditions: hidden, iceberg, pegged}

\begin{definition}[Hidden order, iceberg order]\label{def:mx:order-types-and-their-uses:hidden}
A \emph{hidden order}\index{hidden order} rests in the book and can execute against incoming orders but appears in no market-data feed. An \emph{iceberg
order}\index{iceberg order} (reserve order) displays only a slice of its size; when the slice is executed, the venue displays the next slice from the
reserve, at the back of the queue.
\end{definition}

A \omterm{def:mx:order-types-and-their-uses:hidden}{hidden order} gives up two things for its secrecy. At its price it ranks behind every displayed order, whatever their arrival times (the priority
\texttt{firm.lob} implements), and, being invisible, it attracts no one: a displayed bid tells sellers where to sell, a hidden one waits to be found.
An iceberg's refreshed slice is a new arrival and joins the back of the queue with a new reference on the feed (chapter~1), so a large iceberg trades
one slice per pass of the queue.

\begin{definition}[Pegged order, midpoint peg]\label{def:mx:order-types-and-their-uses:peg}
A \emph{pegged order}\index{pegged order} has no fixed price: the venue sets it from a reference quote and resets it whenever the reference moves,
within an optional limit. A primary peg follows the best quote on its own side; a \emph{midpoint peg}\index{midpoint peg} rests hidden at the midpoint
of the best bid and best ask.
\end{definition}

A \omterm{def:mx:order-types-and-their-uses:peg}{midpoint peg} trades at half the spread: a buyer and a seller who both peg to the midpoint meet there, each saving half a spread against crossing it.
That is the business of dark pools (chapter~9), and the reason exchanges offer \omterm{def:mx:order-types-and-their-uses:peg}{midpoint pegs} too. The rule text of one exchange shows how much fine
print the order type carries.

\begin{dated}{2026-06}{Pegging and minimum quantity on Nasdaq}\label{dat:mx:order-types-and-their-uses:nasdaq}
Nasdaq's Rule 4703, as reproduced in a rule filing of June 2026: an order with Pegging receives a new timestamp whenever its price is updated (it is
evaluated as a newly entered order); an order with Midpoint Pegging is removed when the inside bid and offer become crossed or missing, re-entered at
the new midpoint, and cancelled if no valid price returns within one second; midpoint-pegged orders are cancelled when a trading halt is declared;
pegging orders are subject to a collar and are cancelled where they would execute more than USD~0.25 or 5\% (whichever is greater) worse than the
national best bid and offer; an order with a Minimum Quantity attribute must be at least a round lot and may not be displayed.
\end{dated}

\begin{omfigure}
\begin{tikzpicture}[font=\footnotesize,x=1cm,y=1cm]
\node[anchor=west,font=\small\bfseries] at (-0.2,3.3) {iceberg};
\foreach \x/\w/\l in {0/1.3/A,1.35/0.9/B} {\draw[fill=omDef!18,draw=omDef] (\x,2.4) rectangle ++(\w,0.55); \node at (\x+\w/2,2.675) {\l};}
\draw[fill=omExo!40,draw=omExo,thick] (2.3,2.4) rectangle ++(0.95,0.55); \node at (2.775,2.675) {slice};
\draw[fill=gray!12,draw=gray,dashed] (2.3,1.6) rectangle ++(2.5,0.55); \node at (3.55,1.875) {hidden reserve};
\draw[->,thick,omExo] (3.3,2.675) to[out=0,in=180] (5.0,2.675);
\draw[fill=omExo!40,draw=omExo,thick] (5.05,2.4) rectangle ++(0.95,0.55); \node at (5.525,2.675) {next};
\node[anchor=west,align=left] at (6.1,2.675) {when the slice is executed, the next one\\ joins the back, under a new reference};
\node[anchor=west,font=\small\bfseries] at (-0.2,0.8) {midpoint peg};
\draw[thick,omDef] (0,0) -- (3,0); \node[anchor=east] at (-0.05,0) {bid 100.00};
\draw[thick,omThm] (0,0.5) -- (3,0.5); \node[anchor=east] at (-0.05,0.5) {ask 100.02};
\draw[dashed,gray] (0,0.25) -- (3,0.25);
\draw[fill=gray!20,draw=gray] (1.1,0.13) rectangle ++(0.8,0.24);
\node[anchor=west,align=left] at (3.3,0.25) {rests hidden at the mid, 100.01; moves when either quote moves;\\ trades at half the spread against either side};
\end{tikzpicture}

\omcaption{Two display conditions. Top: an iceberg shows one slice; when it is consumed, the venue shows the next slice at the back of the queue.
Bottom: a \omterm{def:mx:order-types-and-their-uses:peg}{midpoint peg} rests, invisible, halfway between the best bid and ask.}\label{fig:mx:order-types-and-their-uses:display}
\end{omfigure}

An iceberg hides its size but not its behaviour. On \texttt{firm.exchsim}, with Book~7's simulated flow as background, an agent sells 30\,000 shares
at the best ask showing 300 at a time, first as a venue iceberg, then as an algorithm that sends a new 300-share order itself after each slice fills. A
detector reads only the public feed: after an execution removes a displayed order entirely, it looks for an add at the same side and price within 50
milliseconds and of exactly the removed order's size. \Cref{fig:mx:order-types-and-their-uses:iceberg} shows the arms race. The venue's refresh arrives
in the same event as the execution and is caught every time (99 of 99), with 20 false alarms in the half hour from other orders that happened to
match. An algorithm that refills at once is caught 62\% of the time (its refills that met a market that had moved executed on arrival and rested
smaller); refilling after a random delay averaging 200 milliseconds drops that to 18\%, and randomising the size by up to 30\% to 6\%. Loosening the
detector (one second, sizes within 35\%) brings detection back to 47\% but with 378 false alarms for 34 detections. Randomised delays and sizes are
why execution algorithms that refill their own orders are harder to see than venue icebergs.

\begin{omfigure}
\begin{tikzpicture}
\begin{axis}[name=rate,width=7.2cm,height=5.2cm,title={\small refills detected},ybar,bar width=7pt,ymin=0,ymax=1.1,
  xtick={0,1,2,3},xticklabels={venue,algo,algo 200 ms,algo 200 ms $\pm$30\%},x tick label style={font=\footnotesize,rotate=20,anchor=east},
  tick label style={font=\footnotesize},grid=major,grid style={gray!25},enlarge x limits=0.18,
  legend columns=2,legend style={font=\footnotesize,at={(1.12,-0.42)},anchor=north,draw=none,fill=none,
  /tikz/every even column/.append style={column sep=8pt}},table/col sep=comma]
\addplot[fill=omDef!60,draw=omDef] table[x=k,y=strict]{figdata/microstructure/02-order-types-and-their-uses/iceberg.csv};
\addplot[fill=omThm!60,draw=omThm] table[x=k,y=loose]{figdata/microstructure/02-order-types-and-their-uses/iceberg.csv};
\legend{strict detector (50 ms; exact size),loose detector (1 s; size within 35\%)}
\end{axis}
\begin{axis}[at={(rate.east)},anchor=west,xshift=1.2cm,width=7.2cm,height=5.2cm,title={\small false alarms in half an hour},ybar,bar width=7pt,
  ymin=0,ymax=420,xtick={0,1,2,3},xticklabels={venue,algo,algo 200 ms,algo 200 ms $\pm$30\%},
  x tick label style={font=\footnotesize,rotate=20,anchor=east},tick label style={font=\footnotesize},grid=major,grid style={gray!25},
  enlarge x limits=0.18,table/col sep=comma]
\addplot[fill=omDef!60,draw=omDef] table[x=k,y=strict_fa]{figdata/microstructure/02-order-types-and-their-uses/iceberg.csv};
\addplot[fill=omThm!60,draw=omThm] table[x=k,y=loose_fa]{figdata/microstructure/02-order-types-and-their-uses/iceberg.csv};
\end{axis}
\end{tikzpicture}

\omcaption{Detecting icebergs from the public feed: the share of true refills a strict and a loose detector flag (left) and their false alarms (right),
for a venue iceberg and for an algorithm refilling immediately, after a random delay (mean 200~ms), and with a random size (up to 30\% either way);
30\,000 shares shown 300 at a time in a simulated half hour. Data: \texttt{mx\_ordertypes.iceberg\_study}.}\label{fig:mx:order-types-and-their-uses:iceberg}
\end{omfigure}

\section{Who uses which, and why}

\begin{center}
\small
\begin{tabular}{@{}p{3.1cm}p{4.9cm}p{5.6cm}@{}}
\toprule
participant & typical orders & why \\
\midrule
retail investor & market, limit, stop & simple; the broker routes (One Quant Book 1, chapter~10) \\
execution algorithm & \omterm{def:mx:the-limit-order-book:orders}{limit orders} at the best, pegs, hidden, iceberg, IOC to dark pools & work a large order without showing it; take liquidity only when needed \\
market maker & post-only limits, mass quotes, IOC to hedge & earn the spread and the rebate; never take by accident \\
latency-sensitive trader & IOC and intermarket sweep orders & take a stale price before it moves, touch nothing else \\
position-protecting trader & stop and stop-limit & get out at a price without watching the screen \\
\bottomrule
\end{tabular}
\end{center}

Order types multiply because each solves one participant's problem inside the priority rules, and each can change who trades first. The 2015
case of the hook is the extreme: order types that re-priced an order to avoid locking another venue's quote, with a priority behind the scenes that
only some members fully knew. The lesson for anyone who trades is to read the rulebook and the technical specification together and test the engine's
behaviour against both; for anyone who builds an engine (chapter~26), to write the rule in one place and derive the documentation from it.

\section{A limit order is a free option}

\begin{definition}[Free option of a limit order]\label{def:mx:order-types-and-their-uses:option}
The \emph{free option of a limit order}\index{free option of a limit order} is the right a resting \omterm{def:mx:the-limit-order-book:orders}{limit order} grants the rest of the market: to buy at
its price (a resting sell is a call written at that strike) or to sell at its price (a resting buy is a put), exercisable by anyone, at any time, until
the order is cancelled.
\end{definition}

Copeland and Galai (1983) described the cost of supplying quotes this way: a dealer's bid and ask are a put and a call written to traders who know more.
The option is exercised when the value moves through the price before the owner can cancel. A European version gives a lower bound on its value.

\begin{proposition}[The value given away by a resting order]\label{prop:mx:order-types-and-their-uses:option}
A sell order rests at distance $d\ge0$ above the efficient price $P^\ast_0$, which moves as a Brownian motion with volatility $\sigma$ per square-root
second. If a trader who knows $P^\ast_T$ buys from it at time $T$ whenever $P^\ast_T$ exceeds its price, the order loses on average
\[
C(d)=\sigma\sqrt T\,\varphi\!\left(\frac{d}{\sigma\sqrt T}\right)-d\left(1-\Phi\!\left(\frac{d}{\sigma\sqrt T}\right)\right).
\]
\end{proposition}

\begin{proof}
The loss is $(P^\ast_T-P^\ast_0-d)^+$ with $P^\ast_T-P^\ast_0\sim\mathcal N(0,\sigma^2T)$: the Bachelier call value (One Quant Book 2, chapter~13).
\end{proof}

\omcode{code/microstructure/02-order-types-and-their-uses/python/mx_ordertypes.py}{202}{209}{The free option of a resting order.}\label{lst:mx:order-types-and-their-uses:option}

In the simulated hour of \texttt{firm.tape} the efficient price moves 0.27 ticks per square-root second. A sell half a tick above value that stays for
10 seconds gives away 0.15 ticks per share in expectation, and 0.60 ticks if it stays a minute: an order's value decays with the time it stays
unattended, which is why liquidity providers cancel so much (chapter~1). The measured cost is larger, because the informed can exercise at the best
moment, not at a fixed one: marked to the efficient price 10 seconds after each trade, the resting side of trades against informed aggressors lost 1.22
ticks a share; against uninformed ones it gained 0.58; informed flow was 9.2\% of volume, and all resting orders together gained 0.41. Chapters 4 and 5
turn this arithmetic into the theory of the spread.

\section{Tutorial: one book, every order type}\label{tut:mx:order-types-and-their-uses:tut}

\begin{tutorial}
\textbf{Goal.} See the order types act on one engine, provoke a stop cascade, and try to find an iceberg in the public feed. \textbf{End state:}
the table of section~2, \cref{fig:mx:order-types-and-their-uses:cascade,fig:mx:order-types-and-their-uses:iceberg}.
\begin{tutsteps}
\item \textbf{Validation.} Read the order-type rules that run before matching:
\omcode{code/firm/ordertypes/firm_ordertypes.py}{71}{92}{Validation of an entered order in \texttt{firm.ordertypes}.}\label{lst:mx:order-types-and-their-uses:validate}
\item \textbf{Showcase.} \texttt{mx\_ordertypes.showcase()} sends the eight orders to a bare engine and returns the reports of each.
\item \textbf{Cascade.} \texttt{stop\_cascade(n)} for 0 to 20 stops, with \omterm{def:mx:order-types-and-their-uses:stop}{stop orders} and with \omterm{def:mx:order-types-and-their-uses:stop}{stop-limit orders}.
\item \textbf{Icebergs.} \texttt{iceberg\_study(mode, delay\_ms, jitter)} for the four cases, with the strict and the loose detector; draw with
\texttt{fig\_ordertypes.py}.
\item \textbf{The option.} \texttt{passive\_markout(session())} and \texttt{free\_option(0.5, sigma\_per\_sqrt\_second(session()), 10)}.
\end{tutsteps}
\textbf{What to change next.} Space the stops two cents apart and find the spacing at which the cascade stops by itself; give the detector the
iceberg's price as a prior (it knows the level) and see how many false alarms that removes.
\end{tutorial}

\section{Build: the order-type layer}\label{bld:mx:order-types-and-their-uses:ordertypes}

\begin{build}
\bfield{Purpose} One place for the rules each order type adds to matching, used by the exchange simulator's engine in three languages.
\bfield{Interface} \texttt{EOrder} (a \texttt{firm\_lob.Order} with owner, firm, \omterm{def:mx:order-types-and-their-uses:tif}{time in force}, display, post-only, display quantity, reserve, minimum
quantity, self-trade prevention group and mode, stop price, limit, container); \texttt{validate(msg, inst, phase, frozen, band)} returning a rejection
reason; \texttt{slice\_for\_display}; \texttt{peg\_target(o, book, pegged)}; \texttt{stp\_conflict}, \texttt{stp\_actions(mode)};
\texttt{stop\_triggered(o, last)}; \texttt{marketable(side, limit, price)}.
\bfield{Rules} Validation reasons are the order-entry protocol's rejection codes; a \omterm{def:mx:order-types-and-their-uses:peg}{midpoint peg} is priced from the displayed quotes and capped by its
limit; a primary peg follows the best unpegged displayed price on its side; minimum quantity only on IOC orders and \omterm{def:mx:order-types-and-their-uses:peg}{midpoint pegs}; a display quantity
only on visible orders; stops trigger on trades, not quotes.
\bfield{Acceptance tests} \texttt{code/firm/exchsim/tests/test\_engine\_rules.py}: post-only, IOC and FOK, \omterm{def:mx:the-limit-order-book:orders}{market orders} and bands, icebergs, self-trade
prevention in each mode, replace priority, \omterm{def:mx:order-types-and-their-uses:peg}{midpoint pegs} with a minimum, stops, the opening auction, throttles and duplicates; the C++20 and Rust engines
reproduce the Python engine byte for byte on the shared journal (chapter~26).
\bfield{Stretch} Discretionary orders (displayed at one price, willing to trade at another); a price-sliding order that re-prices instead of locking an
away market, with its priority stated.
\end{build}

\begin{omsources}
\item SEC, \emph{In the Matter of EDGA Exchange, Inc. and EDGX Exchange, Inc.}, Release No.~74032, 12~January 2015.
\item Nasdaq Rule 4703 (Order Attributes), as reproduced in SR-NASDAQ-2026-055, Release No.~34-105718, June 2026.
\item T.~E.~Copeland and D.~Galai, ``Information effects on the bid-ask spread'', \emph{Journal of Finance} 38(5), 1983.
\end{omsources}

\section{Exercises}

\begin{exercise}[$\star$]\label{exo:mx:order-types-and-their-uses:1}
The best ask is 100.01 for 300 shares and 100.02 for 500. What does a buy limit at 100.02 for 600 shares do, and what does it leave in the book?
\end{exercise}

\begin{exercise}[$\star$]\label{exo:mx:order-types-and-their-uses:2}
With 300 shares offered at 100.01 and nothing behind them within the limit, what do a fill-or-kill and an immediate-or-cancel buy of 400 at 100.01 each do?
\end{exercise}

\begin{exercise}[$\star$]\label{exo:mx:order-types-and-their-uses:3}
Why does a \omterm{def:mx:order-types-and-their-uses:hidden}{hidden order} rank behind a displayed order at the same price even if it arrived first?
\end{exercise}

\begin{exercise}[$\star\star$]\label{exo:mx:order-types-and-their-uses:4}
An iceberg of 10\,000 shares shows 500. If it is executed in full, how many slices appear on the feed, and how many times does it go to the back of
the queue?
\end{exercise}

\begin{exercise}[$\star\star$]\label{exo:mx:order-types-and-their-uses:5}
A quote rests one tick from value; the efficient price moves 0.5 ticks per square-root second. What does the proposition give for a quote left for
30 seconds?
\end{exercise}

\begin{exercise}[$\star\star$]\label{exo:mx:order-types-and-their-uses:6}
Explain why \omterm{def:mx:order-types-and-their-uses:stop}{stop-limit orders} shorten the cascade of \cref{fig:mx:order-types-and-their-uses:cascade}, and what their owners give up.
\end{exercise}

\begin{exercise}[$\star\star\star$]\label{exo:mx:order-types-and-their-uses:7}
\emph{Coding.} Make the iceberg agent choose its refill price as the current best ask instead of its first price, rerun the strict detector, and explain
the change.
\end{exercise}

\begin{exercise}[$\star\star\star$]\label{exo:mx:order-types-and-their-uses:8}
\emph{Find the flaw.} ``\omterm{def:mx:order-types-and-their-uses:peg}{Midpoint pegs} save half the spread on every trade, so an execution algorithm should send all its orders as \omterm{def:mx:order-types-and-their-uses:peg}{midpoint pegs}.''
\end{exercise}

\section{Problem: The Iceberg in the Tape}

\begin{problem}[{Weekend problem --- the iceberg in the tape}]\label{pb:mx:order-types-and-their-uses:1}
A fund sells 30\,000 shares at the ask, 300 shown at a time. A competitor watches the public feed. You have \texttt{firm.exchsim}, Book~7's background
flow and the chapter's detector.

\medskip\noindent\textbf{Part I --- The order types.}
\begin{enumerate}
\item What is the difference between an iceberg and a \omterm{def:mx:order-types-and-their-uses:hidden}{hidden order} in priority and in what the feed shows?
\item What happens to an iceberg's priority when its slice is refreshed, and why?
\item What does a \omterm{def:mx:order-types-and-their-uses:tif}{fill-or-kill order} protect its sender from, and what does it cost?
\item Why do \omterm{def:mx:order-types-and-their-uses:stop}{stop orders} make a falling book fall faster?
\end{enumerate}

\medskip\noindent\textbf{Part II --- The cascade.}
\begin{enumerate}[resume]
\item How far does a 3\,000-share market sell move the price with no stops, and with twenty?
\item How many shares trade in the twenty-stop case?
\item How far does the price fall with \omterm{def:mx:order-types-and-their-uses:stop}{stop-limit orders} one cent below their triggers?
\item What would you change in the bid ladder to stop the cascade?
\end{enumerate}

\medskip\noindent\textbf{Part III --- The detector.}
\begin{enumerate}[resume]
\item Describe the strict detector.
\item How many refills does the venue iceberg make, and how many does the strict detector catch? With how many false alarms?
\item Why does an algorithm that refills at once get caught less often than the venue's iceberg?
\item What do a 200-millisecond random delay and a 30\% random size do to detection?
\end{enumerate}

\medskip\noindent\textbf{Part IV --- The race.}
\begin{enumerate}[resume]
\item What does the loose detector detect with the delayed and randomised refills, and at what cost in false alarms?
\item State the \emph{named result}: the detection and false-alarm rates of the refill detector against the refill's delay and size randomisation.
\item What further information could the watcher use to separate true refills from coincidences?
\item What does the seller lose by delaying refills?
\item Is detecting icebergs from public data legal? Is trading on the detection?
\item Which order type would you use to sell 30\,000 shares without any refill pattern?
\item How would a venue design an iceberg that is harder to detect?
\item In one sentence: what does an order type reveal about its owner?
\end{enumerate}
\end{problem}

\section{Interview questions}

\begin{interviewq}[$\star$ \iqroles{trader}]\label{iq:mx:order-types-and-their-uses:1}
When would you use an \omterm{def:mx:order-types-and-their-uses:tif}{immediate-or-cancel order} rather than a \omterm{def:mx:the-limit-order-book:orders}{limit order} that rests?
\end{interviewq}

\begin{interviewq}[$\star$ \iqroles{trader, researcher}]\label{iq:mx:order-types-and-their-uses:2}
What is a stop-loss order, and why can it execute far from its stop price?
\end{interviewq}

\begin{interviewq}[$\star\star$ \iqroles{researcher}]\label{iq:mx:order-types-and-their-uses:3}
Why is a resting \omterm{def:mx:the-limit-order-book:orders}{limit order} like a written option? Who holds the option, and what is its strike?
\end{interviewq}

\begin{interviewq}[$\star\star$ \iqroles{developer}]\label{iq:mx:order-types-and-their-uses:4}
How would you implement an \omterm{def:mx:order-types-and-their-uses:hidden}{iceberg order} in a \omterm{def:mx:the-limit-order-book:engine}{matching engine}? What goes on the feed when the displayed slice is executed?
\end{interviewq}

\begin{interviewq}[$\star\star$ \iqroles{trader, developer}]\label{iq:mx:order-types-and-their-uses:5}
A midpoint-pegged order is resting when the market locks (bid equals ask). What should the engine do, and why?
\end{interviewq}

\begin{interviewq}[$\star\star\star$ \iqroles{researcher, trader}]\label{iq:mx:order-types-and-their-uses:6}
How would you detect a large hidden seller from public market data alone?
\end{interviewq}
